\chapter[Identificación de métodos builders]{Enfoques para la identificación de Métodos Builders}
\label{cap:builders}

El análisis de software es una tarea crucial en el campo de la Ingeniería de Software, ya sea para generar tests para el software bajo test (SUT)  o para realizar verificación de modelos de software. Estos enfoques requieren que los programadores identifiquen manualmente un subconjunto de los métodos de un módulo con el fin de dirigir el análisis. En general, al analizar un módulo, los programadores seleccionan un subconjunto de sus métodos que serán considerados como constructores de objetos (\emph{builders}) para definir lo que se conoce como un controlador o "driver", que se utilizará para construir objetos automáticamente para el análisis, combinándolos de manera no determinista, aleatoria, etc. Esto requiere una inspección minuciosa del módulo y su API, ya que la exhaustividad relativa del análisis (omitir métodos importantes puede evitar sistemáticamente la generación de diferentes objetos) y su eficiencia (las diferentes combinaciones acotadas de métodos crecen de manera exponencial a medida que aumenta el número de métodos) se ven afectados por la selección.

En este capitulo, proponemos tres enfoques para seleccionar automáticamente un conjunto de constructores de objetos a partir de la API de un módulo, basados en un algoritmo evolutivo, un algoritmo greedy y un algoritmo basado en clases de equivalencias. Estos algoritmos favorecen conjuntos de métodos cuyas combinaciones conducen a la generación de conjuntos más grandes de objetos o a aquellos conjuntos que obtienen mejor cobertura de ramas del SUT. Además, los algoritmos tienen en cuenta otras características de estos conjuntos de métodos, tratando de priorizar la selección de aquellos con menor cantidad de métodos y con más simples parámetros.

% Para evaluar experimentalmente nuestra propuesta, hemos realizado una evaluación en un conjunto de clases con estado de referencia, que representan casos de uso típicos. Los resultados muestran que nuestro enfoque puede identificar automáticamente conjuntos de constructores de objetos que son suficientes (se pueden utilizar para crear cualquier instancia del módulo) y mínimos (no contienen métodos superfluos), en un tiempo razonable.

% En resumen, en este trabajo abordamos el desafío de seleccionar automáticamente un conjunto óptimo de constructores de objetos a partir de la API de un módulo, con el objetivo de mejorar tanto la exhaustividad como la eficiencia del análisis de software. Nuestros resultados experimentales demuestran la viabilidad y efectividad de nuestra propuesta, lo que abre nuevas perspectivas para la automatización de la selección de constructores de objetos en el análisis de software.


\section{Motivacion}
\label{sec:motivacion}
\begin{table}[H]
\center
{\scriptsize
\begin{tabular}{|l|l|l|l|}
\hline
No. &Return type & Method name & Obs? \\
\hline
    0 && NCL() & no \\
    1& & NCL(int) & no \\
    2&& NCL(Collection) & no \\
    3&boolean & add(Object) & no \\
    4&void&add(int,Object) & no \\
    5&boolean&addAll(Collection) & no\\
    6&boolean&addAll(int,Collection) & no \\
    7&boolean&addFirst(Object) & no \\
    8&boolean&addLast(Object) & no\\
    9&void&clear() & no\\
    10&boolean&contains(Object) & yes \\
    11&boolean&containsAll(Collection) & yes \\
    12&boolean&equals(Object) & yes \\
    13&Object&get(int) & yes\\
    14&Object&getFirst() &yes \\
    15&Object&getLast() & yes\\
    16&int&indexOf(Object) &yes\\
    17&boolean&isEmpty() & yes\\
    18&Iterator&iterator() & no\\
    19&int&lastIndexOf(Object) &yes \\
    20&ListIterator&listIterator() &no \\
    21&ListIterator&listIterator(int) & no\\
    22&Object&remove(int) &no\\
    23&boolean&remove(Object) & no \\
    24&boolean&removeAll(Collection) & no \\
    25&Object&removeFirst() &no\\
    26&Object&removeLast() &no\\
    27&boolean&retainAll(Collection) &no \\
    28&Object&set(int,Object) &no\\
    29&int&size() &yes\\
    30&List&subList(int,int) & no \\
    31&Object[]&toArray() & yes \\
    32&Object[]&toArray(Object[]) &yes\\
    33&String&toString() & yes \\
\hline
\end{tabular}
}
\caption{Apache's NodeCachingLinkedList API}
\label{tab:ncl-api}
\end{table} 
En esta sección, motivamos nuestro enfoque mediante un ejemplo práctico. La estructura de datos NodeCachingLinkedList (NCL) de Apache \cite{apache} consta de una lista principal circular doblemente enlazada  que contiene los elementos de la colección y una lista secundaria simplemente enlazada que actúa como caché para los nodos que se han eliminado de la lista principal. Los nodos almacenados en la caché pueden ser reutilizados y añadidos de nuevo a la lista principal al insertar elementos en ella. Gracias a su caché, en las aplicaciones en las que las inserciones y eliminaciones de la lista son muy frecuentes, NCL puede reducir significativamente la sobrecarga necesaria para la asignación de memoria y la recolección de basura de los nodos. 
Como ilustración, la Figura~\ref{fig:ncl-instances} muestra las tres instancias de NCL que se pueden construir con exactamente dos nodos.


\begin{figure}[H]
    \centering
    \includegraphics[width=0.85\textwidth]{NCL-instances.png}
    \caption{Three NodeCachingLinkedList instances with exactly two nodes}
    \label{fig:ncl-instances}
\end{figure}


NCL tiene una API muy completa, como se muestra en la Tabla~\ref{tab:ncl-api}. Sin embargo, para construir cualquier objeto de NCL factible, sólo se necesitan algunos métodos de la API. Por ejemplo, las combinaciones de los métodos en la Figura~\ref{fig:NCLbuilders}, cuando se instancian con los parámetros apropiados, se pueden utilizar para construir cualquier objeto de NCL acotado con el scope deseado. Por lo tanto, los siguientes métodos:
\\
\begin{lstlisting}[numbers=none,label=fig:NCLbuilders, caption=Conjunto de metodos sufiente para NCL]
  (0)  NodeCachingLinkedList()
  (7)  addFirst(Object)
  (25) removeFirst()
\end{lstlisting}

 son un ejemplo de un conjunto suficiente de constructores (\emph{builders}) de NCL. Hay que tener en cuenta que, después de utilizar el método constructor, la lista principal de NCL se puede rellenar simplemente utilizando el método \texttt{addFirst}. Sin embargo, si queremos generar instancias en las que la lista secundaria (lista caché) no esté vacía, podemos hacerlo a través del método \texttt{removeFirst}, como sugiere el conjunto suficiente de constructores. Para la mayoría de las técnicas de análisis automáticas, nos gustaría considerar tantos escenarios (entradas) variables como sea posible, de ahí la motivación para construir estas con un un conjunto suficientes de métodos de la API. Además, los metodos \emph{builders} de la figura~\ref{fig:NCLbuilders} también son mínimos, ya que la falta de alguno de ellos implicaría que algunos objetos de NCL ya no se pueden construir bajo el scope deseado.

Hay que tener en cuenta que puede haber muchos conjuntos de constructores suficientes. Por ejemplo, se pueden obtener constructores suficientes reemplazando el método \texttt{addFirst} en la Figura~\ref{fig:NCLbuilders} por cualquier otra variante de \texttt{add} que se muestra en la Figura~\ref{fig:NCLadds}, ya que para cualquier manera de llenar la lista principal de NCL con \texttt{addFirst}, existe una forma diferente de construir el mismo objeto utilizando otra variante de \texttt{add} (quizás invocada con diferentes parámetros y cambiando el orden de ejecución).
\\
\begin{lstlisting}[numbers=none,label=fig:NCLadds, caption=Variantes del método 'Add' que puedo ser utilizado para rellanar la lista principal en NCL, captionpos=b, frame=tb , basicstyle=\scriptsize]
  (3) add(Object)
  (4) add(int,Object)
  (7) addFirst(Object)
  (8) addLast(Object)
\end{lstlisting}


Además, es importante remarcar la importancia de obtener subconjuntos que sean mínimos, lo que significa que no contengan métodos adicionales. Esto es crucial para utilizar la combinación de estos métodos en la construcción de objetos. Cuantos más métodos tenga el subconjunto de builders, más costosa será la combinación de los mismos para generar los objetos.

Los siguientes subconjuntos son suficientes pero no son mínimos para NCL:
\\

\begin{lstlisting}[numbers=none,label=fig:NCLadds, frame=tb , basicstyle=\scriptsize]
  (0)  NodeCachingLinkedList()
  (7)  addFirst(Object)
  (3)  addF(Object)
  (25) removeFirst()
\end{lstlisting}

\begin{lstlisting}[numbers=none,label=fig:NCLadds, caption= Conjuntos de metodos builders suficientes pero no mínimos, captionpos=b, frame=tb , basicstyle=\scriptsize]
  (0)  NodeCachingLinkedList()
  (7)  addFirst(Object)
  (3)  remove(Object)
  (25) removeFirst()
\end{lstlisting}

También observamos que cuanto más simples sean los parámetros de una rutina, más fácil será utilizarla para generar entradas en el contexto de un análisis de programas. Por ejemplo, entre las alternativas de rutinas para añadir dato a la lista principal de NCL (Figura.~\ref{fig:NCLadds}), \texttt{add(int,Object)} recibe más parámetros que los otros tres métodos, por lo tanto es más difícil generar parámetros para brindale a esta rutina a la hora de generar entradas. Esto hace que las otras tres alternativas tengan mas preferencia sobre esta. Así, nuestro enfoque tiene en cuenta el número de parámetros y sus complejidades para seleccionar los constructores mejores posibles.

Muchos métodos en la tabla \ref{tab:ncl-api} están marcados como observadores (columna Obs?), lo que significa que no modifican los objetos en los que operan, ni son útiles para crear objetos no primitivos. Por lo tanto, los observadores siempre son superfluos y nunca deben incluirse en un conjunto de constructores mínimos. Nuestro enfoque intenta reconocerlos de antemano y descartarlos de la búsqueda para reducir significativamente el espacio de búsqueda.

Todos estos problemas planteados en este capítulo fueron lo que nos motivó a desarrollar enfoques para obtener automáticamente métodos constructores de manera eficiente y mínima. Para lograrlo, utilizamos una variedad de técnicas y algoritmos que explicaremos en la próxima sección. 

\section{Función de Valoración}
\label{sec:fitness}
En un algoritmo genético, la función de valoración, también conocida como función de aptitud o función de fitness, es esencial para evaluar y medir la calidad de las soluciones individuales en una población \cite{goldberg1989genetic}. Esta función asigna un valor numérico a cada individuo en función de su grado de idoneidad para resolver el problema en cuestión. El objetivo principal de la función de valoración es proporcionar una medida cuantitativa de la calidad de cada solución candidata en relación con los criterios específicos del problema.

La función de valoración se diseña de acuerdo con la naturaleza del problema y los objetivos específicos del algoritmo genético. Por ejemplo, si estamos resolviendo un problema de optimización en el que buscamos maximizar una función objetivo, la función de valoración puede asignar un valor más alto a las soluciones que se acercan más a la solución óptima. Por otro lado, si estamos resolviendo un problema de minimización, la función de valoración puede asignar un valor más alto a las soluciones que se alejan más de la solución óptima.

En el caso de algoritmos greedy, la función de valoración juega un papel similar pero con un enfoque más local. En lugar de trabajar con una población de soluciones, los algoritmos greedy toman decisiones en cada paso basándose en una evaluación local de las opciones disponibles. La función de valoración en un algoritmo greedy se utiliza para seleccionar la mejor opción en cada iteración o paso del algoritmo, maximizando o minimizando la función de valoración según el objetivo del problema \cite{cormen2009introduction}.

Por lo tanto, tanto en algoritmos genéticos como en algoritmos greedy, la función de valoración desempeña un papel crucial para evaluar y comparar soluciones candidatas en función de criterios específicos del problema. Proporciona una medida cuantitativa de la calidad de las soluciones y guía el proceso de selección y toma de decisiones en cada etapa del algoritmo. El diseño adecuado de la función de valoración es esencial para obtener resultados óptimos y eficientes en la resolución de problemas.
A continuacion, explicaremos cada funcion de valoracion implementada en nuestros algoritmos.







\subsection{Fitness: Generador Exhaustivo}

Dado un candidato que representa un conjunto de métodos $M$, nuestra función de valoración intenta calcular una aproximación del número de objetos acotados que se pueden construir utilizando combinaciones de métodos habilitados en el cromosoma. Los candidatos con valores de aptitud más altos se estima que construyen más objetos que aquellos que tienen valores de aptitud más pequeños.
Nuestra función de valoración tiene como objetivo principal maximizar la cantidad de objetos que puedo generar utilizando un Generador exhaustivo (explicado en la sección de Preliminares \ref{sec:beapiIntro}. Este enfoque nos permite evaluar la capacidad de los candidatos para construir una variedad de objetos y proporciona una base para la selección y mejora de los mejores candidatos en nuestro algoritmo.

Esta función de valoración en nuestra tesis se utiliza para evaluar la calidad de un candidato, el cual representa un conjunto de métodos $M$. Para ello, nuestra función devuelve un valor real que se construye de la siguiente manera:

\[
f(M) = \text{{\#Objectos}}(M) \ , \ (\text{{\#MétodosAPI}} - \text{{\#M}}) + w_1 \times \text{{PP}}(M)
\]

La función de valoración se compone de tres componentes. En primer lugar, $\text{{\#Objetos}}(M)$ representa el número de objetos generados por el conjunto de métodos $M$. La coma en la fórmula indica que este valor es independiente de los otros componentes, y se le da prioridad a la cantidad de objetos creados con el subconjunto de métodos $M$.

El segundo componente, $(\text{{\#MétodosAPI}} - \text{{\#M}})$, refleja la diferencia entre el número total de métodos disponibles en la API y el número de métodos en el conjunto $M$. Este término se utiliza para desempatar en caso de que otro subconjunto de métodos $M_1$ construya la misma cantidad de objetos que $M$.

Finalmente, $\text{{PP}}(M)$ es un factor de ponderación que permite ajustar la importancia de tener menor complejidad y cantidad de parámetros en el conjunto de métodos $M$. Este valor es entre 1 y 0 es solo para ajustar en caso de que haya empate de cantidad de objetos y métodos.
\cacho{Explicar el peso de los parametros}
Esta función de valoración nos proporciona una medida cuantitativa de la calidad del candidato representado por el conjunto de métodos $M$. Los candidatos con valores de función de valoración más altos se consideran de mayor calidad y son preferidos en los algoritmos genéticos y en los algoritmos greedys que explicaremos en la sección \ref{sec:algorithms}.
Idealmente, nos gustaría explorar todos los objetos factibles dentro de un límite pequeño $k$ que se pueden construir utilizando los métodos del cromosoma actual. En otras palabras, necesitamos un generador exhaustivo acotado para el conjunto de métodos $BE(M, k)$. El límite $k$ representa el número máximo de objetos que se pueden crear para cada clase (en la Figura \ref{fig:ncl-instances}, el número de nodos en los objetos NCL está acotado por $k=2$) y el número máximo de valores primitivos disponibles (por ejemplo, enteros del 0 a $k-1$).

Para este propósito, desarrollamos la herramienta BEAPI, que se discute con más detalle en la Capítulo \ref{cap:beapi}. En resumen, primero exploramos exhaustivamente todas las posibles combinaciones de secuencias de los métodos de $M$. Luego, utilizamos un conjunto fijo de valores primitivos (enteros del 0 a $k-1$) con los cuales probar nuestros métodos cuando requieren valores primitivos.

En segundo lugar, descartamos las secuencias de métodos que crean objetos con más de $k$ objetos (de cualquier tipo) para evitar construir objetos más grandes de lo necesario. Para lograr esto, canonicamos los objetos generados por la ejecución de cada secuencia y descartamos la secuencia si algún objeto tiene un índice igual o mayor que $k$.

En tercer lugar, ampliamos esta generación con coincidencia de estado. Esto se debe a que, en la generación de pruebas, a menudo hay muchas secuencias de pruebas que producen el mismo objeto. Por ejemplo, insertar en una colección y luego eliminar el mismo elemento resulta en muchos casos en exactamente la misma estructura antes de la inserción. Nuestro enfoque asume que las ejecuciones de rutinas son deterministas con respecto a sus entradas. Bajo esta suposición, se deduce que, para generar un conjunto exhaustivo acotado de estructuras, solo necesitamos guardar una secuencia de prueba para crear cada estructura diferente en el conjunto, y que todas las siguientes secuencias de prueba que generen la misma estructura se pueden descartar.

Nuestra justificación para usar conjuntos acotados de objetos es similar a la \emph{hipótesis de la cota pequeña} \cite{Andoni:2003}.Si un conjunto de métodos no puede usarse para construir objetos pequeños que permitan diferenciarlo de otro conjunto de métodos, es poco probable que estos dos conjuntos puedan distinguirse con objetos más grandes. Esta hipótesis se mantuvo durante nuestra evaluación empírica en todos nuestros casos de estudio.

Para obtener más información sobre el algoritmo y más detalles sobre BEAPI, invitamos al lector a consultar el Capítulo \ref{cap:beapi}.


\subsection{Fitness: Cobertura de ramas con Randoop}

En el contexto de nuestra tesis, hemos realizado modificaciones en la herramienta \emph{Randoop} descrita en la sección \ref{sec:feedback-directed-test-gen}.

Nuestras modificaciones se centran en dar prioridad a un subconjunto específico de métodos, representado en la sección anterior por el conjunto $M$. El objetivo de esta modificación es generar secuencias de pruebas que se enfoquen principalmente en la ejecución de estos métodos seleccionados. Sin embargo, es importante destacar que esto no implica excluir por completo otros métodos, ya que aún buscamos ejercitar la API en su totalidad para obtener una buena cobertura de todos los métodos.

La motivación detrás de esta modificación radica en que los subconjuntos de métodos que obtienen una mayor cobertura de la API nos proporcionarán indicios valiosos sobre la calidad de dicho subconjunto. Por lo tanto, nuestra función de valoración tiene como objetivo maximizar la cobertura de ramas de la API bajo prueba, tal como se describe en la sección de preliminares \ref{sec:coverage}.

La función de valoración que utilizamos para evaluar la calidad de un candidato, que representa un conjunto de métodos $M$, se construye de la siguiente manera:

\[
f(M) = \text{{\#CoberturaRama}}(M) \, (\text{{\#MétodosAPI}} - \text{{\#M}}) + w_1 \times \text{{PP}}(M)
\]

La función de valoración consta de tres componentes principales. En primer lugar, $\text{{\#CoberturaRama}}(M)$ representa el número de ramas cubiertas por las suites de prueba generadas por Randoop utilizando el subconjunto de métodos $M$. Este componente refleja la cantidad de ramas cubiertas específicamente por este subconjunto priorizado de métodos.

El segundo componente, $\text{{\#MétodosAPI}} - \text{{\#M}}$, tiene en cuenta la cantidad de métodos restantes en la API que no están presentes en el subconjunto $M$. Este componente incentiva a que el conjunto de métodos sea menor.

Por último, penalizamos sobre los parámetros de cada método en el subconjunto como lo hacíamos en la sección anterior.

En resumen, nuestras modificaciones en Randoop nos permiten priorizar un subconjunto de métodos durante la generación de secuencias de pruebas, y nuestra función de valoración se encarga de maximizar la cobertura de ramas de la API. Estas adaptaciones nos proporcionan una mejor comprensión de la calidad y efectividad de los subconjuntos de métodos generados.


\section{Algoritmos}


\label{sec:algorithms}
Para encontrar un conjunto suficiente de constructores a partir de una API de un programa, diseñamos tres algoritmos de búsqueda que se describen a continuación. El primero es una modificación de los algoritmos genéticos, el segundo es un algoritmo Greedy más precisamente una variante de un algoritmo \emph{Hill Climbing}, y el tercero es un algoritmo de búsqueda de acuerdo a las clases de equivalencias de los subconjuntos respecto al valor de su función de valoración.
Antes de ejecutar cualquiera de nuestros algoritmos, aplicamos una técnica de análisis estático para eliminar de la búsqueda aquellos métodos que son irrelevantes para nuestros métodos builders. En otras palabras, necesitamos identificar y eliminar aquellos métodos que son considerados observadores en la API.

Los métodos observadores son aquellos que no modifican el estado de un objeto al ejecutarse. En el ejemplo motivador presentado en la sección \ref{sec:motivation} y la tabla \ref{tab:ncl-api}, se puede observar que existen muchos métodos clasificados como observadores. Esto nos indica que estos métodos no son relevantes para ser considerados como métodos builders.
Para realizar esta identificación de métodos observadores, utilizamos una herramienta de análisis estático llamada \emph{Infer}\footnote{https://fbinfer.com/}. Esta herramienta ha sido desarrollada por \emph{Facebook} y ofrece varias funcionalidades específicas. En el contexto de nuestra tesis, utilizamos únicamente su analizador estático para detectar qué métodos son considerados puros u observadores.

El análisis estático realizado con \emph{Infer} nos permite obtener una lista de métodos que son clasificados como observadores en la API bajo prueba. Esta información es utilizada como paso previo a la ejecución de nuestros algoritmos, para asegurarnos de que solo consideramos los métodos relevantes y adecuados para comenzar la búsqueda de los métodos \emph{builders}.

En resumen, el análisis estático realizado con \emph{Infer} nos permite eliminar de la búsqueda los métodos observadores, aquellos que no modifican el estado de los objetos, garantizando así que nos enfocamos en los métodos más relevantes para nuestros métodos builders.

\subsection{Algoritmo Genético}
\label{alg:approachGA}
En esta sección presentamos los detalles para la detección de los subconjuntos de métodos builders utilizando un algoritmo evolutivo. Para lograr realizar esto, implementamos un algoritmo genético (\ref{sec:geneticoPrev}) que busca el subconjunto de métodos que sean mínimo y suficientes, que describiremos a continuación.


% % basados en una estrategia de escalada de colinas \cite{Russell:2009}.
% Durante esta sección explicaré en detalle cada algoritmo.


% Dare mas detalles de este algoritmo en la seccion {TODO!!!!}.


\subsubsection{Cromosomas}
\label{ge:cromosomas}

Los elementos de la población de nuestro algoritmo genético son subconjuntos de métodos de la API. Para representar esto, necesitamos codificarlos como \emph{cromosomas}. En nuestro caso, necesitamos alguna forma de representar los métodos de la clase como un vector. Por ejemplo, consideremos el ejemplo de NCL explicado en el capítulo anterior, que tiene 34 métodos (consultar Tabla \ref{tab:ncl-api}). Para representar nuestras posibles soluciones, crearemos un vector de genes booleanos, donde cada posición $i$ es verdadera si y solo si el cromosoma contiene el $i$-ésimo método de la API. Para hacer esto, enumeraremos los métodos de la API desde 0 hasta $n$, donde $n$ es el número total de métodos en la API. Cabe destacar que cada cromosoma tendrá la misma longitud, lo cual es común en los algoritmos genéticos.

Para el ejemplo de NCL, que tiene 34 métodos (Tabla \ref{tab:ncl-api}), tendríamos la siguiente representación del cromosoma:

\begin{center}
$c = \begin{array}{@{}cccccc@{}}
\scriptsize g_0 & \scriptsize g_1 & \scriptsize g_2 & \ldots & \scriptsize g_{34}
\end{array}$
\end{center}

Aquí, cada gen (\emph{g}) es un valor booleano que indica la presencia o ausencia del método en ese cromosoma.

Como ejemplo más concreto, consideremos la representación del cromosoma que representa el conjunto de métodos suficientes y mínimos para NCL, como se muestra en la Figura \ref{fig:NCLbuilders}:
\begin{center}
$c = \begin{array}{|*{8}{c|}}
\hline
0 & 1 & \ldots & 7 & \ldots & 25 & \ldots & 33 \\
\hline
1 & 0 & \ldots & 1 & \ldots & 1 & \ldots & 0 \\
\hline
\end{array}$
\end{center}

En este caso, las posiciones 0, 7 y 25 están establecidas como verdaderas (según el orden asignado a los métodos de la API), mientras que las demás posiciones están establecidas como falsas.


\subsubsection{Poblacion Inicial}

El proceso comienza con un conjunto de individuos llamado Población. Recordemos que cada individuo es una solución al problema que deseas resolver.

En nuestro algoritmo genético, la población inicial se crea generando aleatoriamente un conjunto de cromosomas, donde cada cromosoma representa una solución potencial al problema. Cada gen en el cromosoma se establece aleatoriamente como 1 o 0, lo que indica la presencia o ausencia del método que se codificó en esa posición. Además, todos nuestros cromosomas iniciales tienen un tamaño igual a la cantidad de métodos que tiene la API bajo test.

Es importante destacar que el número de cromosomas en la población inicial desempeña un papel fundamental en el algoritmo genético, ya que proporciona la base para la evolución y mejora gradual de las soluciones a lo largo de las generaciones. En nuestro algoritmo, tenemos una población inicial de 100 individuos. A medida que el algoritmo avanza, se aplican operadores genéticos como la selección, el cruce y la mutación para crear nuevas soluciones a partir de la población actual. A continuación, se explican estos operadores en detalle.

\subsubsection{Operadores Geneticos}

A continuación explicaremos los operadores genéticos principales utilizados en el algoritmo evolutivo. Cada uno desempeña un papel importante en la exploración y explotación del espacio de búsqueda y en la mejora de la calidad de las soluciones a lo largo de las generaciones.

\paragraph{Cross-over (Cruce)}

El cruzamiento (o \emph{cross-over} en inglés) es un operador genético que combina una parte de dos cromosomas $c1$ y $c2$ elegidos aleatoriamente para crear un nuevo cromosoma descendiente. Esa parte que toma de cada cromosoma padre depende de la configuración con la que se lo setee. En nuestro algoritmo, utilizamos un ratio de 0.30. Además, un cruzamiento de dos puntos de cruce dentro de los cromosomas de los padres. Su objetivo principal es explorar y explotar la información genética existente en la población.
El cruce de varios puntos permite una mayor variedad en los descendientes generados, ya que se intercambian segmentos más largos de genes entre los padres. Esto puede ayudar a explorar el espacio de búsqueda de soluciones de manera más efectiva y encontrar soluciones potencialmente mejores.

\paragraph{Mutacion}

La mutación es un operador fundamental en un algoritmo genético que introduce pequeños cambios aleatorios en los cromosomas para explorar nuevas soluciones y mantener la diversidad dentro de la población. Juega un papel crucial en la prevención de la convergencia prematura y garantiza que el algoritmo explore de manera efectiva el espacio de búsqueda. Existen varios operadores de mutación que se pueden utilizar en un algoritmo genético. Uno de los operadores de mutación más comúnmente utilizados es la Mutación por Inversión de Bits, que es adecuada para algoritmos genéticos codificados en binario como en el caso de nuestra codificación.
En la Mutación por Inversión de Bits, se selecciona uno o más bits aleatorios en el cromosoma y se invierten. Si un bit está inicialmente establecido en 1, se cambia a 0, y viceversa. Esta alteración aleatoria introduce cambios sutiles en la información del cromosoma, creando potencialmente nuevas soluciones que no estaban presentes en la población original. El número de bits invertidos y sus posiciones se determinan típicamente mediante una baja probabilidad de mutación.


\paragraph{Selección}
La operación de selección es un componente clave en nuestro algoritmo genético, ya que determina qué individuos serán preservados en la siguiente generación. En nuestro enfoque, hemos desarrollado un operador de selección que se basa en el nivel de aptitud de cada individuo.

Utilizamos un enfoque de selección tipo torneo (\emph{tournament}) con un tamaño de torneo de 4 en nuestro algoritmo. En este enfoque, se selecciona aleatoriamente un grupo de individuos de la población y se compite entre ellos. El individuo con el mejor valor de fitness en el torneo es seleccionado para formar parte de la próxima generación. Este proceso se repite hasta que se hayan seleccionado todos los individuos necesarios.

Es importante destacar que el operador de selección basado en torneo tiene varias ventajas. No solo permite seleccionar a los individuos más aptos, sino que también introduce una presión de selección ajustable. Esto significa que los individuos más débiles tienen menos posibilidades de ser seleccionados en torneos más grandes, lo que ayuda a mejorar gradualmente la calidad de la población a lo largo de las generaciones.

\cacho{Algortimo? esta en preliminares, es standard para Genetico.}
%Pongo ejemplo de codigo???%



\subsection{Hill Climbing}
\label{alg:approachHC}
El algoritmo Hill Climbing, también conocido como búsqueda por ascenso de colina\cite{Russell:2009}, es un algoritmo de búsqueda local que se utiliza para optimizar una función de valoracion. El objetivo del algoritmo es encontrar la solución óptima para un problema determinado.  Su enfoque se basa en realizar movimientos ascendentes, es decir, buscar soluciones que mejoren continuamente el valor objetivo o la función de evaluación.
Este algoritmo es una algoritmo Greedy(Perezoso) explicado en la seccion de preliminares, (\ref{sec:greedyPrev}). Este selecciona un buen estado sucesor para el momento actual, sin pensar en dónde ir a continuación. 
La idea principal detrás del algoritmo Hill Climbing es comenzar con una solución inicial y, en cada iteración, realizar un movimiento hacia una solución vecina que mejore el valor objetivo. Este proceso se repite hasta que no se pueda encontrar una solución vecina que mejore el valor actual. En ese punto, el algoritmo se detiene y devuelve la mejor solución encontrada hasta el momento.
Es importante destacar que el algoritmo Hill Climbing puede quedar atrapado en óptimos locales, es decir, puede converger hacia soluciones que son mejores en comparación con sus vecinos inmediatos, pero no son óptimas en el contexto global

La representación que hemos utilizado del problema fue igual a la que utilizamos para representar cromosomas en el Algoritmo genético explicado en la sección anterior. Esto quiere decir que utilizamos vectores de 0 y 1 para representar una posible solución. 

A continuación mostraremos un pseudocódigo del algoritmo \emph{Hill Climbing} que hemos implementado:

\begin{algorithm}[H]
  \caption{Algoritmo de Hill Climbing}
  \label{algo:hill_climbing}
  \SetAlgoLined
  \KwResult{Solución óptima $curr$}
  $curr \gets c$\; 
  
  \While{existe un mejor candidato}{
    $S \gets$ GenerarSucesores($curr$)\;
    
    $best \gets$ SeleccionarMejorSucesor($S$)\;
    
    \If{$f(best) > f(curr)$}{
      $curr \gets best$\;
    }
  }
  \Return{$best$}\;
\end{algorithm}

Este algoritmo representa el esquema básico de Hill Climbing, comienza calculando la funcion de valoración de todos los singletones ${c}$ de métodos constructores.  El mejor de los singletones (mayor objectos puedo crear con ese constructor) se establece como el candidato actual $curr$, y Hill Climbing inicia un proceso de búsqueda típico e iterativo (Línea 1).

En cada iteración (Línea 2 a 8), \emph{Hill Climbing} calcula $f(succ)$ para cada succ $\in$ $GenerarSucesores(curr)$. El método que genera los sucesores nos devuelve un conjunto de posibles soluciones que se crean a partir de sumarle un método a la solución óptima corriente ($curr$). Es decir, los sucesores  $S$ generados con \emph{GenerarSucesores(curr)} de un candidato $curr$ son los conjuntos {$curr\cup{mi}$}, para cada $mi$ $\in$ API (Línea 3).
Sea $best$ el sucesor con el valor de valoración más alto. Observe que $best$ tiene exactamente un método más que el mejor candidato de la iteración anterior, $curr$ (Línea 4).

Si $f(best) > f(curr)$, los métodos en $best$ se pueden utilizar para crear un conjunto más grande de estructuras que los en $curr$ (Línea 5 a 7). Por lo tanto, \emph{Hill Climbing} asigna $best$ a $cur$r y continúa con la siguiente iteración. En cambio, si $f(best) <= f(curr)$, $curr$ ya genera el conjunto más grande de estructuras posible (no se puede agregar ningún método que aumente el número de estructuras generadas a partir de $curr$). En este punto, $curr$ se devuelve como el conjunto de constructores identificados. (Linea 9) 
\cacho{
En este algoritmo no se necesitan setear ningun parametros para guiar la busqueda a que sea mas efectiva}

Se puede observar que este algoritmo puede quedar atrapado en un máximo local y no generar alguna combinación específica que podría ser aún mejor. Aquí se puede ver que es un algoritmo greedy, ya que obtiene una solución rápida pero puede no ser eficiente en encontrar la mejor solución global. 


\subsection{Clases de equivalencia}
\label{alg:approachCE}
 Las clases de equivalencia son una técnica utilizada para agrupar conjuntos de datos de entrada en categorías o clases que tienen un comportamiento similar o producen resultados equivalentes. Esta técnica es ampliamente utilizada en el diseño y la realización de pruebas de software.
 En términos generales, una clase de equivalencia representa  un conjunto de que se espera que produzcan resultados idénticos. La idea es que si un conjunto de la clase de equivalencia produce un resultado entonces todas las demás entradas de esa clase deberían producir el mismo resultado.
 Al trabajar con clases de equivalencia, se selecciona una entrada representativa, llamada caso de prueba, de cada clase para ser evaluada. En lugar de probar todas las posibles entradas, se eligen casos de prueba que representen cada clase de equivalencia para minimizar la cantidad total de pruebas necesarias.

 Bajo esta introducción, hemos desarrollado un tercer algoritmo donde agrupamos en clases de equivalencia aquellos subconjuntos de métodos que tengan el mismo valor de valoración con alguna de las funciones explicadas en \ref{sec:fitness}.


\begin{algorithm}[H]
  \caption{Algoritmo basado en Clases de Equivalencia}
  \label{algo:clases_equivalencia}
  \SetAlgoLined
  \KwResult{Conjunto de metodos builders $best$}
  $curr \gets c$\; 
  $equivalenceClasses \gets$ CrearClasesDeEquivalencia($curr$)\;
  
  \While{se ha creado una nueva clase de equivalencia}{
    $newCandidates \gets$ CandidatosPorClase($equivalenceClasses$)\;
    
    \ForEach{$candidate$ en $newCandidates$}{
      $successors \gets$ GenerarSucesores($candidate$)\;   
      
      \ForEach{$successor$ en $successors$}{
        $key \gets$ f($successor$)\;
        $equivalenceClasses[key]$.put($successor$)\;
      }
    }
        % $best \gets$ SeleccionarMejorCandidato($candidates$)\;
  }
  
  $result \gets$ obtenerMejor($best$)\;
  \Return{$result$}\;
\end{algorithm}

En este algoritmo, se comienza obteniendo los conjuntos singletons (Línea 1), que son los métodos constructores individuales de la misma manera que el algoritmo de \ref{algo:hill_climbing}. A partir de estos singletons, se selecciona el mejor candidato inicial $curr$, que servirá como punto de partida. Luego, se crean las clases de equivalencia basadas en la función de valoración y agrupando los candidatos que tienen el mismo valor de valoración. En el algortimo las clases de equivlencias estan guardadas en $equivalenceClasses$ (Línea 2).

A continuación, se itera (Línea 5 a 12) mientras se haya creado una nueva clase de equivalencia. En cada iteración, se generan nuevos candidatos por cada clase de equivalencia. Es decir, se selecciona un candidato de cada clase de equivalencia (el de menor métodos y menor parámetros) $CandidatosPorClase(equivalenceClasses)$ y se lo guarda en un cola ($newCandidates$) (Línea 5). A continuación, se generan los sucesores (línea 6) para cada candidato elegido, $GenerarSucesores(candidate)$, utilizando el mismo enfoque que en el algoritmo de \emph{Hill Climbing}. Estos se guardan en $successors$ Es decir, si tenemos $N$ clases de equivalencias vamos a tener $N$ candidatos a los cuales le vamos a calcular sus sucesores.  Luego, por cada uno de estos sucesores en  se le calcula su función de valoración y de acuerdo a esta se lo guarda en la clase de equivalencia representada por el valor de valoración.
El algoritmo itera hasta que no hayamos creado una nueva clase de equivalencia. Esto quiere decir que no hay cambios y hemos explorado todas las alternativas posibles. 
Luego vamos a obtener el mejor conjunto de métodos builders, accediendo a la clase de equivalencia con mayor key y obteniendo el representante de ese conjunto que tenga menor métodos y parámetros. Esto lo realiza el método $obtenerMejor(equivalenceClasses)$

